\documentclass{article}

\usepackage[utf8]{inputenc}
\usepackage{amsmath, esint}
\usepackage{wasysym}
\usepackage{qrcode}
\usepackage[colorlinks]{hyperref}
\usepackage{lmodern}
\usepackage{graphicx}
\usepackage{xcolor}
\usepackage[left=2cm, top=3cm, right=2cm]{geometry}
\usepackage{minted}
\usepackage{booktabs}
\usepackage{svg}
\usepackage{xcolor}
\definecolor{LightGray}{gray}{0.975}
\hypersetup{
  linkcolor=blue,
  urlcolor=blue
}

\title{ML101 \\ Homework 03: A `gentle' Introduction to Classification and Model Selection.}
\author{Andrés Oswaldo Calderón Romero, Ph.D.}
\date{\today}

\begin{document}

\maketitle

\section{Introduction}
\label{sec:intro}

This assignment introduces fundamental concepts in supervised classification and linear model selection. The primary objectives are:

\begin{itemize}
    \item \textbf{Non-Parametric Classification:} Build intuitive foundations for the $K$-Nearest Neighbors ($K$-NN) algorithm through selected introductory literature.
    \item \textbf{Applied Classification:} Implement foundational classification algorithms using R/Python, covering Sections 4.7.1--4.7.5 of \textit{An Introduction to Statistical Learning} to support upcoming laboratory work.
    \item \textbf{Model Selection and Regularization:} Prepare for upcoming coursework via a flipped-classroom study of Chapter~6 video lectures, focusing on subset selection, stepwise methods, and validation-based test error estimation.
\end{itemize}

\section{Reading Material}
\label{sec:reading}

To develop an intuitive understanding of the $K$-Nearest Neighbors ($K$-NN) algorithm, select and review \textbf{one} of the following introductory articles:

\begin{itemize}
  \item \textit{``K-Nearest Neighbor (KNN) Explained''} by Diego Lopez Yse (2023). Available at \href{https://www.pinecone.io/learn/k-nearest-neighbor/}{Pinecone.com}.
  \item \textit{``What is the k-nearest neighbors (KNN) algorithm?''} by Eda Kavlakoglu (2022). Available at \href{https://www.ibm.com/think/topics/knn}{IBM.com}.
  \item \textit{``K-Nearest Neighbor (KNN) Algorithm''} by GeeksforGeeks Community (2026). Available at \href{https://www.geeksforgeeks.org/machine-learning/k-nearest-neighbours/}{GeeksforGeeks.org}.
\end{itemize}

This reading is intended for conceptual preparation; no deliverable is required for this section.

\section{Lab: Classification Methods}
\label{sec:lab}

To gain hands-on experience with classification methods, work through the Chapter~4 laboratory exercises from \textit{An Introduction to Statistical Learning Using R} (Sections 4.7.1--4.7.5, pp.~171--181). Given the fact that all of you are pythonistas, you can cover the similar sections in the Python book.

While no formal deliverable is required for this section, working through these exercises is essential preparation for the upcoming laboratory assignment.

\section{Flipped Classroom}
\label{sec:flipped}

To prepare for next week's lectures, review the assigned video recordings for Chapter~6 (``Linear Model Selection and Regularization''). These lectures are provided by Stanford Online and are hosted on their YouTube channel and the edX platform\footnote{Hastie, T., \& Tibshirani, R. (2023). \textit{Statistical learning} (STATSX0001) [MOOC]. edX; Stanford Online. \url{https://www.statlearning.com/online-courses}}.

Refer to Table~\ref{tab:videos} for the required playlist. You must take handwritten notes for each lecture video, which will constitute the primary deliverable for this assignment.

\begin{table}
  \centering
  \caption{Required Video Lectures for Chapter 6: Linear Model Selection and Regularization}
  \label{tab:videos}
  \begin{tabular}{c l c c c c c}
    \toprule
      Section & Title & Duration & Links & & & \\
    \midrule
      6.1 & Introduction and Best Subset Selection & 13:45 &
      \href{https://youtu.be/nsv5rEV3mVI?si=l85FAQRbUhYeuaQ0}{YouTube} &
      \href{https://drive.google.com/drive/folders/1nqH4_lnJA7RIghWuMf0VKDTElWQlDkc3?usp=sharing}{Video} &
      \href{https://drive.google.com/drive/folders/1nqH4_lnJA7RIghWuMf0VKDTElWQlDkc3?usp=sharing}{Subtitles} &
      \href{https://drive.google.com/drive/folders/1nqH4_lnJA7RIghWuMf0VKDTElWQlDkc3?usp=sharing}{Transcript} \\
      6.2 & Stepwise Selection & 12:27 &
      \href{https://youtu.be/ynXq-Gw1xfE?si=MTcilvJexgyJULhq}{YouTube} &
      \href{https://drive.google.com/drive/folders/1nqH4_lnJA7RIghWuMf0VKDTElWQlDkc3?usp=sharing}{Video} &
      \href{https://drive.google.com/drive/folders/1nqH4_lnJA7RIghWuMf0VKDTElWQlDkc3?usp=sharing}{Subtitles} &
      \href{https://drive.google.com/drive/folders/1nqH4_lnJA7RIghWuMf0VKDTElWQlDkc3?usp=sharing}{Transcript} \\
      6.3 & Backward Stepwise Selection & 5:27 &
      \href{https://youtu.be/c5aI9cowjRI?si=G60SfaXBt6DU2FxT}{YouTube} &
      \href{https://drive.google.com/drive/folders/1nqH4_lnJA7RIghWuMf0VKDTElWQlDkc3?usp=sharing}{Video} &
      \href{https://drive.google.com/drive/folders/1nqH4_lnJA7RIghWuMf0VKDTElWQlDkc3?usp=sharing}{Subtitles} &
      \href{https://drive.google.com/drive/folders/1nqH4_lnJA7RIghWuMf0VKDTElWQlDkc3?usp=sharing}{Transcript} \\
      6.4 & Estimating test error & 14:07 &
      \href{https://youtu.be/48P-oV6cH44?si=OMwV3cCBfeSiSZl7}{YouTube} &
      \href{https://drive.google.com/drive/folders/1nqH4_lnJA7RIghWuMf0VKDTElWQlDkc3?usp=sharing}{Video} &
      \href{https://drive.google.com/drive/folders/1nqH4_lnJA7RIghWuMf0VKDTElWQlDkc3?usp=sharing}{Subtitles} &
      \href{https://drive.google.com/drive/folders/1nqH4_lnJA7RIghWuMf0VKDTElWQlDkc3?usp=sharing}{Transcript} \\
      6.5 & Validation and cross validation & 8:44 &
      \href{https://youtu.be/mzb5Xs58bb0?si=RNIM797vRbxlgyex}{YouTube} &
      \href{https://drive.google.com/drive/folders/1nqH4_lnJA7RIghWuMf0VKDTElWQlDkc3?usp=sharing}{Video} &
      \href{https://drive.google.com/drive/folders/1nqH4_lnJA7RIghWuMf0VKDTElWQlDkc3?usp=sharing}{Subtitles} &
      \href{https://drive.google.com/drive/folders/1nqH4_lnJA7RIghWuMf0VKDTElWQlDkc3?usp=sharing}{Transcript} \\
    \bottomrule
  \end{tabular}
\end{table}

\section{Expected Deliverable}
\label{sec:deliverables}

You will prepare a simple report sharing your handwritten lecture notes of the videos presented in section~\ref{sec:flipped}.  Please submit your report in \textbf{PDF format} through the corresponding BrightSpace\texttrademark{} assignment for next week.

\vspace{5mm}
Happy Learning! \includesvg[width=4mm]{figures/sunglasses}

\end{document}

